\documentclass[12pt,a4paper]{article}
\usepackage[utf8]{vietnam}
\usepackage{amsmath,amssymb}
\usepackage[top=2cm, bottom=2cm, left=2cm, right=2cm]{geometry}
\usepackage{ex_test} % Load the package with new designs

\begin{document}

\section{DEMO CÁC MÔI TRƯỜNG CAO CẤP}
\subsection{Môi trường Lý thuyết (Định nghĩa, Định lý, Tính chất)}

\begin{dn}{Số Nguyên Tố}{}
    Số nguyên tố là số tự nhiên lớn hơn 1 chỉ có hai ước là 1 và chính nó.
\end{dn}

\begin{dl}{Định lý Pytago}{}
    Trong một tam giác vuông, bình phương cạnh huyền bằng tổng bình phương hai cạnh góc vuông:
    \[ a^2 = b^2 + c^2 \]
\end{dl}

\begin{tc}{Tính chất trọng tâm}{}
    Trọng tâm của tam giác chia mỗi đường trung tuyến thành hai đoạn thẳng theo tỉ lệ 2:1, kể từ đỉnh.
\end{tc}

\subsection{Môi trường Bài tập & Ví dụ}

\begin{vd}{Giải phương trình bậc hai}{}
    Giải phương trình: $x^2 - 3x + 2 = 0$.
    
    \textit{Lời giải:}
    Ta có $\Delta = (-3)^2 - 4(1)(2) = 1 > 0$.
    Phương trình có hai nghiệm phân biệt:
    \[ x_1 = \frac{3-1}{2} = 1; \quad x_2 = \frac{3+1}{2} = 2. \]
\end{vd}

\begin{bt}{Tính tích phân}{}
    Tính tích phân sau:
    \[ I = \int_0^1 x e^x \mathrm{d}x \]
\end{bt}

\subsection{Môi trường Đặc biệt}

\begin{dang}[Dạng 1: Phương trình mũ cơ bản]
    Phương pháp giải: Đưa về cùng cơ số hoặc đặt ẩn phụ.
    \[ a^{f(x)} = a^{g(x)} \Leftrightarrow f(x) = g(x) \]
\end{dang}

\begin{note}
    \textbf{Lưu ý quan trọng:}
    Khi giải phương trình chứa căn thức, cần đặt điều kiện xác định trước khi bình phương hai vế.
\end{note}

\end{document}
